\uuid{LDHh}
\titre{Ensembles}
\auteur{bodin ; croizat ; sacre}
\organisation{exo7}
\licence{CC BY-NC-SA 4.0}
\datecreate{2022-03-08}
\niveau{L1}
\module{Algèbre}
\chapitre{Logique, ensemble, raisonnement}
\sousChapitre{Ensemble}
\difficulte{1}
\theme{Ensembles}
\exo7qcmid{hilisit.logique.025}

\contenu{
\begin{question}
\qclassification{100.02}
Soient les ensembles $A = [1,3]$ et $B= \{ 0 , 1 , 2 , 3 \}$. Quelles sont les assertions vraies ?
\begin{answers}
	\bad{$B \subset A$}
	\bad{$A \subset B$}
	\bad{$A \setminus B = ]1,3[$}
	\good{$A \cap B = \{ 1,2,3 \}$}
\end{answers}
\begin{explanations}
Il n'y a pas d'inclusion entre $A$ et $B$ : Le nombre $\frac32 \in A $ et $\frac32 \notin B$ d'une part ; et $0 \in B$ avec $0 \notin A$ d'autre part.\\
On a $A \setminus B = ]1,2[ \cup ]2,3[$.\\
En revanche, on a bien $A \cap B = \{1,2,3\}$.
\end{explanations}
\end{question}
}
